\documentclass[a4paper,11pt,oneside]{memoir}

% Castellano
\usepackage[spanish,es-tabla]{babel}
\selectlanguage{spanish}
\usepackage[utf8]{inputenc}
\usepackage{placeins}

\RequirePackage{booktabs}
\RequirePackage[table]{xcolor}
\RequirePackage{xtab}
\RequirePackage{multirow}

% Links
\usepackage[colorlinks]{hyperref}
\hypersetup{
	allcolors = {red}
}

% Ecuaciones
\usepackage{amsmath}

% Rutas de fichero / paquete
\newcommand{\ruta}[1]{{\sffamily #1}}

% Párrafos
\nonzeroparskip


% Imagenes
\usepackage{graphicx}
\newcommand{\imagen}[2]{
	\begin{figure}[!h]
		\centering
		\includegraphics[width=0.9\textwidth]{#1}
		\caption{#2}\label{fig:#1}
	\end{figure}
	\FloatBarrier
}

\newcommand{\imagenflotante}[2]{
	\begin{figure}%[!h]
		\centering
		\includegraphics[width=0.9\textwidth]{#1}
		\caption{#2}\label{fig:#1}
	\end{figure}
}



% El comando \figura nos permite insertar figuras comodamente, y utilizando
% siempre el mismo formato. Los parametros son:
% 1 -> Porcentaje del ancho de página que ocupará la figura (de 0 a 1)
% 2 --> Fichero de la imagen
% 3 --> Texto a pie de imagen
% 4 --> Etiqueta (label) para referencias
% 5 --> Opciones que queramos pasarle al \includegraphics
% 6 --> Opciones de posicionamiento a pasarle a \begin{figure}
\newcommand{\figuraConPosicion}[6]{%
  \setlength{\anchoFloat}{#1\textwidth}%
  \addtolength{\anchoFloat}{-4\fboxsep}%
  \setlength{\anchoFigura}{\anchoFloat}%
  \begin{figure}[#6]
    \begin{center}%
      \Ovalbox{%
        \begin{minipage}{\anchoFloat}%
          \begin{center}%
            \includegraphics[width=\anchoFigura,#5]{#2}%
            \caption{#3}%
            \label{#4}%
          \end{center}%
        \end{minipage}
      }%
    \end{center}%
  \end{figure}%
}

%
% Comando para incluir imágenes en formato apaisado (sin marco).
\newcommand{\figuraApaisadaSinMarco}[5]{%
  \begin{figure}%
    \begin{center}%
    \includegraphics[angle=90,height=#1\textheight,#5]{#2}%
    \caption{#3}%
    \label{#4}%
    \end{center}%
  \end{figure}%
}
% Para las tablas
\newcommand{\otoprule}{\midrule [\heavyrulewidth]}
%
% Nuevo comando para tablas pequeñas (menos de una página).
\newcommand{\tablaSmall}[5]{%
 \begin{table}
  \begin{center}
   \rowcolors {2}{gray!35}{}
   \begin{tabular}{#2}
    \toprule
    #4
    \otoprule
    #5
    \bottomrule
   \end{tabular}
   \caption{#1}
   \label{tabla:#3}
  \end{center}
 \end{table}
}

%
% Nuevo comando para tablas pequeñas (menos de una página).
\newcommand{\tablaSmallSinColores}[5]{%
 \begin{table}[H]
  \begin{center}
   \begin{tabular}{#2}
    \toprule
    #4
    \otoprule
    #5
    \bottomrule
   \end{tabular}
   \caption{#1}
   \label{tabla:#3}
  \end{center}
 \end{table}
}

\newcommand{\tablaApaisadaSmall}[5]{%
\begin{landscape}
  \begin{table}
   \begin{center}
    \rowcolors {2}{gray!35}{}
    \begin{tabular}{#2}
     \toprule
     #4
     \otoprule
     #5
     \bottomrule
    \end{tabular}
    \caption{#1}
    \label{tabla:#3}
   \end{center}
  \end{table}
\end{landscape}
}

%
% Nuevo comando para tablas grandes con cabecera y filas alternas coloreadas en gris.
\newcommand{\tabla}[6]{%
  \begin{center}
    \tablefirsthead{
      \toprule
      #5
      \otoprule
    }
    \tablehead{
      \multicolumn{#3}{l}{\small\sl continúa desde la página anterior}\\
      \toprule
      #5
      \otoprule
    }
    \tabletail{
      \hline
      \multicolumn{#3}{r}{\small\sl continúa en la página siguiente}\\
    }
    \tablelasttail{
      \hline
    }
    \bottomcaption{#1}
    \rowcolors {2}{gray!35}{}
    \begin{xtabular}{#2}
      #6
      \bottomrule
    \end{xtabular}
    \label{tabla:#4}
  \end{center}
}

%
% Nuevo comando para tablas grandes con cabecera.
\newcommand{\tablaSinColores}[6]{%
  \begin{center}
    \tablefirsthead{
      \toprule
      #5
      \otoprule
    }
    \tablehead{
      \multicolumn{#3}{l}{\small\sl continúa desde la página anterior}\\
      \toprule
      #5
      \otoprule
    }
    \tabletail{
      \hline
      \multicolumn{#3}{r}{\small\sl continúa en la página siguiente}\\
    }
    \tablelasttail{
      \hline
    }
    \bottomcaption{#1}
    \begin{xtabular}{#2}
      #6
      \bottomrule
    \end{xtabular}
    \label{tabla:#4}
  \end{center}
}

%
% Nuevo comando para tablas grandes sin cabecera.
\newcommand{\tablaSinCabecera}[5]{%
  \begin{center}
    \tablefirsthead{
      \toprule
    }
    \tablehead{
      \multicolumn{#3}{l}{\small\sl continúa desde la página anterior}\\
      \hline
    }
    \tabletail{
      \hline
      \multicolumn{#3}{r}{\small\sl continúa en la página siguiente}\\
    }
    \tablelasttail{
      \hline
    }
    \bottomcaption{#1}
  \begin{xtabular}{#2}
    #5
   \bottomrule
  \end{xtabular}
  \label{tabla:#4}
  \end{center}
}



\definecolor{cgoLight}{HTML}{EEEEEE}
\definecolor{cgoExtralight}{HTML}{FFFFFF}

%
% Nuevo comando para tablas grandes sin cabecera.
\newcommand{\tablaSinCabeceraConBandas}[5]{%
  \begin{center}
    \tablefirsthead{
      \toprule
    }
    \tablehead{
      \multicolumn{#3}{l}{\small\sl continúa desde la página anterior}\\
      \hline
    }
    \tabletail{
      \hline
      \multicolumn{#3}{r}{\small\sl continúa en la página siguiente}\\
    }
    \tablelasttail{
      \hline
    }
    \bottomcaption{#1}
    \rowcolors[]{1}{cgoExtralight}{cgoLight}

  \begin{xtabular}{#2}
    #5
   \bottomrule
  \end{xtabular}
  \label{tabla:#4}
  \end{center}
}


















\graphicspath{ {./img/} }

% Capítulos
\chapterstyle{bianchi}
\newcommand{\capitulo}[2]{
	\setcounter{chapter}{#1}
	\setcounter{section}{0}
	\chapter*{#2}
	\addcontentsline{toc}{chapter}{#2}
	\markboth{#2}{#2}
}

% Apéndices
\renewcommand{\appendixname}{Apéndice}
\renewcommand*\cftappendixname{\appendixname}

\newcommand{\apendice}[1]{
	%\renewcommand{\thechapter}{A}
	\chapter{#1}
}

\renewcommand*\cftappendixname{\appendixname\ }

% Formato de portada
\makeatletter
\usepackage{xcolor}
\newcommand{\tutor}[1]{\def\@tutor{#1}}
\newcommand{\course}[1]{\def\@course{#1}}
\definecolor{cpardoBox}{HTML}{E6E6FF}
\def\maketitle{
  \null
  \thispagestyle{empty}
  % Cabecera ----------------
\noindent\includegraphics[width=\textwidth]{cabecera}\vspace{1cm}%
  \vfill
  % Título proyecto y escudo informática ----------------
  \colorbox{cpardoBox}{%
    \begin{minipage}{.8\textwidth}
      \vspace{.5cm}\Large
      \begin{center}
      \textbf{TFG del Grado en Ingeniería Informática}\vspace{.6cm}\\
      \textbf{\LARGE\@title{}}
      \end{center}
      \vspace{.2cm}
    \end{minipage}

  }%
  \hfill\begin{minipage}{.20\textwidth}
    \includegraphics[width=\textwidth]{escudoInfor}
  \end{minipage}
  \vfill
  % Datos de alumno, curso y tutores ------------------
  \begin{center}%
  {%
    \noindent\LARGE
    Presentado por  \@author{}\\ 
    en Universidad de Burgos ---  \@date{}\\
    Tutor:   \@tutor{Carlos López Nozal}\\
  }%
  \end{center}%
  \null
  \cleardoublepage
  }
\makeatother

\newcommand{\nombre}{Juan Pedro Pascual Vitores} %%% cambio de comando

% Datos de portada
\title{Geoindoor}
\author{\nombre}
\tutor{}
\date{\today}

\begin{document}

\maketitle


\newpage\null\thispagestyle{empty}\newpage


%%%%%%%%%%%%%%%%%%%%%%%%%%%%%%%%%%%%%%%%%%%%%%%%%%%%%%%%%%%%%%%%%%%%%%%%%%%%%%%%%%%%%%%%
\thispagestyle{empty}


\noindent\includegraphics[width=\textwidth]{cabecera}\vspace{1cm}

\noindent D. Carlos López Nozal, profesor del departamento de ingeniería civil , área de lenguaje de sistemas.

\noindent Expone:

\noindent Que el alumno D. \nombre, con DNI 13172380G, ha realizado el Trabajo final de Grado en Ingeniería Informática titulado Geoindoor de TFG. 

\noindent Y que dicho trabajo ha sido realizado por el alumno bajo la dirección del que suscribe, en virtud de lo cual se autoriza su presentación y defensa.

\begin{center} %\large
En Burgos, {\large \today}
\end{center}

\vfill\vfill\vfill

% Author and supervisor
\begin{minipage}{0.45\textwidth}
\begin{flushleft} %\large
Vº. Bº. del Tutor:\\[2cm]
D. Carlos López Nozal
\end{flushleft}
\end{minipage}


\vfill

% para casos con solo un tutor comentar lo anterior
% y descomentar lo siguiente
%Vº. Bº. del Tutor:\\[2cm]
%D. nombre tutor


\newpage\null\thispagestyle{empty}\newpage




\frontmatter

% Abstract en castellano
\renewcommand*\abstractname{Resumen}
\begin{abstract}
La geolocalización es un concepto que se lleva estudiando desde hace años, pero con el avance de la tecnología se ha conseguido acercar la geolocalización al ámbito cotidiano, consiguiendo que llegar a un un determinado lugar nos resulte más fácil.

Habitualmente se ha utilizado la geolocalización para orientarnos en extensiones relativamente amplias. Sin embargo el avance del hardware permite que la geolocalización sea más precisa, con lo que están apareciendo herramientas para que la geolocalización sea aplicada a extensiones de terreno más pequeñas, abriendo un mundo de posibilidades, y precisamente sobre esto trata el proyecto.

Geoindoor es una herramienta con múltiples aplicaciones, que añade servicios de geolocalización dentro edificios. Una aplicación Web permite etiquetar ubicaciones dentro del plano del edificio. Otra aplicación permite consultar las ubicaciones de un edificio , las posiciones en plano del edificio y navegar hasta ellas desde la ubicación actual.
\end{abstract}

\renewcommand*\abstractname{Descriptores}
\begin{abstract}
Geolocalización indoor, sistema de localización, búsqueda de lugares, localización en interiores, rutas predeterminadas. \ldots
\end{abstract}

\clearpage

% Abstract en inglés
\renewcommand*\abstractname{Abstract}
\begin{abstract}
Geolocation is a concept that has been studied since years, but the advancement of technology has been able to bring geolocation closer to us, because of this we arrive to the places faster.

Usually geolocation has been used to orient in large extensions. However, the hardware advancement allows the geolocation to be more accurate, because of this, there are more tools for the geolocation dedicated to smaller terrain extensions, opening a world of possibilities and the project is about it.

Geoindoor is a tool with multiple applications that add geolocation services within buildings. A web application allows you to tag locations within the building plan. Another application that allows you to consult the locations of a building, positions in the building plan and navigate to them from the actual location.
\end{abstract}

\renewcommand*\abstractname{Keywords}
\begin{abstract}
Indoor geolocation, location system, location search, indoor location, predetermined routes. \ldots
\end{abstract}

\clearpage

% Indices
\tableofcontents

\clearpage

\listoffigures

\clearpage

\listoftables
\clearpage

\mainmatter
\capitulo{1}{Introducción}

\texttt{"No se dónde está"}\\
\\
Es habitual escuchar esta frase, que multitud de veces irrumpe en nuestras vidas dificultando nuestras tareas, tareas
tan simples como encontrar un baño, una oficina, o la cafetería. Estos son los típicos problemas que se encuentra uno 
cuando acude por primera vez a un lugar, pero es más grave para aquel que ofrece un servicio o unos productos en un 
determinado lugar y el usuario no lo localiza. Esto puede hacer que el flujo de personas que acudan al establecimiento 
sea mucho menor, lo que provocaría que el servicio estuviera poco solicitado o que los productos ofrecidos no fuesen vistos 
por tantas personas como se desearía. Es aún más grave si hablamos de un negocio el cual necesite que los usuarios localicen
con rapidez donde se ofrecen los servicios o productos. 
 
De ahí la necesidad de la implementación del proyecto, que propone dar servicios de geolocalización indoor. Esto resolvería 
diversos problemas en establecimientos como museos, exposiciones, centros comerciales, hospitales,
o cualquier otro establecimiento de gran tamaño.

La geolocalización indoor además de servir para encontrar localizaciones determinadas, puede ser aplicada para programar 
rutas que se amolden a unos intereses, y así hacer de la experiencia de moverse por nuevos lugares algo más fructífero y gratificante.

\section{Estructura de la memoria}\label{estructura-de-la-memoria}

La memoria tiene la siguiente estructura:

\begin{itemize}
\tightlist
\item
  \textbf{Introducción:} exposición del problema a resolver y como lo resuelve el proyecto.
   Estructura de la memoria y listado de materiales
  adjuntos.
\item
  \textbf{Objetivos del proyecto:} objetivos del proyecto.
\item
  \textbf{Conceptos teóricos:} explicación de los conceptos teóricos necesarios para comprender el proyecto.
\item
  \textbf{Técnicas y herramientas:} técnicas y herramientas utilizadas para la realización de proyecto.
\item
  \textbf{Aspectos relevantes del desarrollo:} explicación de dificultades y obstáculos encontrados en la realización del proyecto.
\item
  \textbf{Trabajos relacionados:} trabajos relacionados con el proyecto.
\item
  \textbf{Conclusiones y líneas de trabajo futuras:} conclusiones
  obtenidas tras realizar el proyecto y expectativas de futuro.
\end{itemize}

Junto a la memoria se proporcionan los siguientes anexos:

\begin{itemize}
\tightlist
\item
  \textbf{Plan del proyecto software:} sección centrada planificación temporal y viabilidad del proyecto.
\item
  \textbf{Especificación de requisitos del software:} sección que se centra en la especificación de requisitos.
\item
  \textbf{Especificación de diseño:} sección que se basa en describir la fase de diseño.
\item
  \textbf{Manual del programador:} sección centrada en la explicación del código fuente (estructura, compilación,
  instalación, ejecución, pruebas,etc... ).
\item
  \textbf{Manual de usuario:} sección centrada en la explicación al usuario de como utilizar la herramienta.
\end{itemize}

\section{Materiales adjuntos}\label{materiales-adjuntos}

Los materiales adjuntados con la memoria son: 

\begin{itemize}
\tightlist
\item
	Architect, parte de la herramienta para crear rutas y edificios.
\item
	Viewer, parte de la herramienta encargada de la visualización .
\item	
	Código fuente de la Rest API encargada de interactuar con la BD.
\item	
	Pruebas de integración.
\item	
	Pruebas de estrés y rendimiento.
\end{itemize}

Recursos en internet:

\begin{itemize}
\tightlist
\item
  Architect (Se lanza en el localhost).
\item
  Viewer (Se lanza en el localhost).
\item
  Rest API https://geoindoorapi.herokuapp.com/ .
\item
  Base de datos de firebase  https://geoindoordb.firebaseio.com/ .
\item
  Repositorio del proyecto .
\end{itemize}

\capitulo{2}{Objetivos del proyecto}

A continuación, se especifican los objetivos que promueven la realización
de este proyecto.
\section{Objetivos generales}\label{objetivos-generales}
\begin{itemize}
\item
Desarrollar una herramienta que permita añadir servicios de geolocalización
en interiores.
\item
Facilitar la localización de lugares dentro de edificios.
\item
Creación de rutas predeterminadas que permitan al usuario tanto diseñar rutas para él mismo,
como para todo el público.
\section{Objetivos técnicos}\label{objetivos-tecnicos}
\item
Uso del API proporcionado por la aplicación \textit{Anyplace}
\item
Uso de Git como sistema de control de versiones junto con la plataforma GitHub. 
\item
Uso de la metodología Scrum.
\item 
Uso de ZenHub como herramienta de gestión de proyectos.
\item
Uso de JS en la parte del servidor (Node.js).
\item
Uso de bases de datos no sql.

\section{Objetivos personales}\label{objetivos-personales}
\item
Entender cómo se realiza un proyecto desde su inicio hasta su final.
\item
Mejorar mis conocimientos y conceptos sobre el desarrollo.
\end{itemize}

\capitulo{3}{Conceptos teóricos}

A continuación se explicarán unos conceptos indispensables para el entendimiento del proyecto.
\section{Geolocalización}\label{geolocalizacion}
La Geolocalización es el concepto más importante en este proyecto, y primero hablaremos un poco del API utilizada y después profundizaremos más en el concepto.

En el proyecto, el mapa que ofrece Google Maps es fundamental, ya que nos permite dibujar sobre el plano y añadir imágenes 
sobre él.

La API utilizada para este proyecto se basa en la API Google Maps, que está escrita en JavaScript \cite{googlemapsjs}.

La clase principal es google.maps.Map, cuyo constructor nos permite crear un mapa dentro de un contenedor HTML. Después encontramos google.maps.LatLng, donde LatLng es un punto en coordenadas geográficas: latitud y longitud.

Para más información se puede visitar la documentación de Google Maps donde aparece información más detallada, ya que no profundizaremos en cuáles, y cómo son las clases y funciones de esta API, sino cómo funciona Google Maps en sí.

Google Maps funciona a través de conexión a internet y la función de GPS del dispositivo, las coordenadas de Google Maps están en el sistema WGS842 \footnote{WGS84 \textit{es un sistema de coordenadas geográficas mundial que permite localizar cualquier punto de la Tierra (sin necesitar otro de referencia) por medio de tres unidades dadas.}} y se mostrará la latitud y la longitud, positiva para Norte y Este, negativa para Sur y Oeste.

Cuando nos referimos a la función GPS, estamos hablando de la función de geolocalización del dispositivo, en la que profundizaremos con más detalle.
\\

\textbf{Definición:} Geolocalización (Geo-localización) está formada por la palabra geo, que tiene origen griego en geos que significa tierra. Y la palabra localización, de las raíces latinas localis  (relativo a lugar), que proviene de locus (lugar).

Por lo tanto, la geolocalización consiste en la ubicación de un ``lugar`` en la ``tierra`` \cite{diccionarioetim}. 

Haciendo una definición más exacta, la geolocalización consiste en ubicar o determinar dónde se encuentra un objeto o lugar en un determinado espacio.

Cuando nos referimos a objeto no tiene porqué ser inanimado, y el determinado espacio al que nos referimos es la tierra como tal, pero no siempre se la pone como plano de referencia. 

Cuando hablamos sobre la geolocalización también debemos hablar de los sistemas de posicionamiento que son los encargados de analizar, manejar y procesar la información geográfica, que nos permite obtener la localización de un objeto u lugar.

El sistema de posicionamiento más conocido en el mundo es el GPS (Global Positioning System ), que es un sistema que permite localizar objetos móviles, gracias a la recepción de señales emitidas por un conjunto de satélites. Para determinar las localizaciones este sistema se sirve de 24 satélites y trilateración. Y es de origen estadounidense, otros sistemas de posicionamiento conocidos son el GLONASS, de origen ruso y Galileo europeo \cite{wikigps}.

La trilateración es un método matemático para determinar las posiciones relativas de los objetos usando la geometría de triángulos y la triangulación. La trilateración para ubicar el objeto usa las posiciones conocidas de dos o más puntos de referencia, y la distancia entre el objeto y cada punto de referencia \cite{wikitrila}.
\\
- Explicación matemática de la trilateración \cite{wikitrila}:
\begin{itemize}
\item 
	Se parte de 3 esferas.
	\imagenResize{0.8}{img/p1_trilateracion}{Paso 1º trilateración }{p1trilateracion}
	\\
\item
	Se halla la x restando la segunda a la primera.
	\imagenResize{0.9}{img/p2_trilateracion}{Paso 2º trilateración }{p2trilateracion}
\item
	Se sustituye la x en la primera esfera, con lo que obtenemos la fórmula de otro círculo.
	\imagenResize{0.8}{img/p3_trilateracion}{Paso 3º trilateración }{p3trilateracion}
\item
	Se iguala la fórmula anterior a la de la tercera esfera, obteniendo.
	\imagenResize{0.8}{img/p4_trilateracion}{Paso 4º trilateración }{p4trilateracion}
\item
	Teniendo las coordenadas del punto solución x e y, despejamos la z de la primera esfera .
	\imagenResize{0.8}{img/p5_trilateracion}{Paso 5º trilateración }{p5trilateracion}
\end{itemize}

De esta manera obtendríamos la solución para los tres puntos x, y y z. Al ser una raíz cuadrada la solución puede ser 0, uno o dos resultados.

\begin{itemize}
\item 
	z = $\sqrt{x}$, $x \prec 0$ : Una de las esferas no intersecciona con las demás (No hay solución real).
\item 
	z = 0 : Las esferas interseccionan justo en un punto.
\item 
	z = $\pm\sqrt{x}$, $x \succ 0$ : Las esferas se interseccionan en 2 puntos.
	
\end{itemize}

Se consigue saber la distancia entre objetos y puntos de referencia gracias a que el dispositivo emite una señal y espera una respuesta de los puntos de referencia. El diferencial de tiempo entre el envío y la recepción se utiliza para determinar la distancia.

Como se muestra en las imágenes en la trilateración en tres dimensiones, los puntos de referencia (satélites) crean una esfera virtual o imaginaria de radio igual a la distancia entre el objeto y el satélite, donde ellos son el epicentro. Se utilizan los puntos de intersección entre esferas para calcular la posición del objeto a localizar.

Los satélites tienen que saber la distancia entre ellos, por lo tanto también mandan señales. Además para hacer un cálculo correcto de distancias se necesitan tener relojes coordinados.
\\
\imagen{img/GPS-Trilateracion}{Explicacion de trilateración.\cite{mejoresappgeo}}



\imagen{img/trilateracion}{Explicacion de trilateración. \cite{animaciontrila}}


La geolocalización por wifi funciona de forma parecida a la de los sistemas de posicionamiento con satélites. Esta vez los puntos de referencia en vez de ser satélites son puntos de acceso wi-fi, y hay varias técnicas y tecnologías para hallar la ubicación de los dispositivos, además de la explicada anteriormente. Una de estas técnicas es la de utilizar RSSI, que indica la intensidad de la señal, a mayor intensidad más cerca está del punto de acceso Wi-Fi que proporciona la señal.  Pero normalmente para la ubicación en interiores se utilizan ondas de radio, ya que se obtiene mayor precisión \cite{wifiloc}. 

\section{Formatos gráficos escalables}\label{formatos-graficos-escalables}

Los formatos gráficos escalables son fundamentales en el proyecto ya que son utilizados a la hora de hacer el dibujado de rutas, así como la creación de iconos.

Cuando hablamos de formatos gráficos escalables, nos estamos refiriendo a aquellos formatos digitales basados en objetos geométricos, cuya labor es hacer que la imagen tenga formas más definidas.
Los formatos más conocidos son SVG, WMF, ODG  y AI.

Debemos entender la importancia de los formatos gráficos escalables en la mejora de imágenes para hacerlas más geométricas.


\subsection{WMF}\label{wmf}

Es un formato de gráficos vectoriales creado por Microsoft, nació a principio de los 90 y ahora está en desuso. Un archivo WMF se utiliza para regenerar una imagen a través Windows GDI \footnote{Graphics Device Interface \textit{interfaz de programación de aplicaciones que se encarga del control gráfico de los dispositivos de salida}} \hypertarget{_Hlk482387012}{. WMF guarda una serie de llamadas a funciones que son enviadas a Windows GDI para regenerar la imagen.} \cite{metaarchivo}

Enhanced Metafile (EMF) es una versión de 32 bits de WMF (16 bits), con comandos adicionales.



\subsection{ODG}\label{ODG}


ODG formato de imagen vectorial enlazado con la versión 2 de OpenDocument de OpenOffice, es un estándar abierto para la creación de dibujos vectoriales, utiliza XML, fue creado por Sun Microsystems y OASIS \cite{definicionodg}.





\subsection{AI}\label{ai}

Formato de imagen vectorial de Adobe Illustrator, está desarrollado por Adobe Systems para la representación de gráficos vectoriales en formato EPS o PDF. El formato AI está compuesto por rutas conectadas mediante puntos, en lugar de datos de imagen. 
En principio el formato AI fue un formato nativo llamado PGF hasta que se compatibilizo con PDF mediante la copia completa de datos PGF en el archivo de formato PDF \cite{extensionai}.




\subsection{SVG}\label{svg}

\begin{flushleft}
Este es el formato gráfico utilizado para el dibujado de rutas y caminos en el proyecto.
\end{flushleft}

SVG (Scalable Vector Graphics) es un formato de imagen vectorial recomendado y desarrollado por W3C (World Wide Web Consortium), es un estándar abierto, trabaja con gráficos vectoriales bidimensionales estáticos y animados, en formato XML.
SVG aparece de forma nativa en multitud de navegadores como Amaya, Mozilla Firefox, Google Chrome, Safari e Internet Explorer. Para navegadores que no tienen SVG de forma nativa existen plugins que permiten mostrar imágenes en formato SVG \cite{scalablevec} \cite{lassoformatos}.
SVG permite tres tipos de objetos graficos:




\begin{itemize}
\item 
	Elementos geométricos vectoriales



\begin{itemize}
\item 
	Rectas 



\item 
	curvas
\end{itemize}


\item 
	Mapa de bits



\item 
	Texto



\end{itemize}


Los objetos gráficos con los que trabaja SVG pueden transformarse y agruparse para después ser compuestos en otros objetos renderizados con anterioridad. 
El dibujado con SVG es dinámico e interactivo, mediante el uso de DOM para SVG (Document Object Model), que mediante ECMAScript \footnote{ECMAScript \textit{es una especificación de lenguaje de programación publicada por ECMA International. Se empezó a desarrollar en 1996 y se basó en JavaScript, fue propuesto como estándar por Netscape Communications Corporation. Actualmente es el estándar ISO 16262.}} o SMIL \footnote{SMIL (Synchronized Multimedia Integration Language)\textit{es un estándar del World Wide Web Consortium (W3C) para presentaciones multimedia, permite integrar audio, video, imágenes, texto o cualquier otro contenido multimedia.}} permite animaciones de gráficos vectoriales sencillas y eficientes
SVG tiene una gran cantidad de manejadores de eventos como ``onMouseOver" y ``onClick", que pueden ser asignados a cualquier objeto SVG. La compatibilidad de SVG con estándares web permite aplicar características como el scripting sobre elementos SVG o XML, de forma simultánea, en la misma página web desde distintos espacios de nombre XML.
Las imágenes SVG pueden ser comprimidas con gzip (SVGZ) \cite{scalablevec}.

\imagen{img/mapabits}{explicación diferencia mapa de bits y gráfico vectorial.\cite{cursografvec}}

\section{Lenguajes de etiquetas}\label{lenguajes-etiuetas}

En cuanto al lenguaje de etiquetas o marcado, al ser una herramienta para la web, la interfaz se nutre de html, con lo cual es una parte fundamental del proyecto.

Los lenguajes de etiquetas, o de marcado, son utilizados para codificar un documento añadiendo etiquetas al texto, las cuales permiten otorgarle unas características o cualidades específicas. En este proyecto es muy importante ya que facilita tanto el procesado de la información, como la muestra de la misma \cite{lenguajemarc}.

No se debe confundir los lenguajes de etiquetas con los de programación, ya que el lenguaje de etiquetas no tiene funciones aritméticas ni variables.

Podemos diferenciar entre 4 tipos de lenguajes de marcado o de etiquetas \cite{lenguajemarcecu} \cite{lmarcas}.
\begin{itemize}
\item
Marcado de presentación.
\item
Marcado de procedimientos.
\item
Marcado descriptivo.
\item
Lenguajes especializados.

\end{itemize}

También existen lenguajes de etiquetado que son difíciles de enmarcar en uno de estos 4 tipos, ya que comparten características de varios de ellos, como podría ser el caso de HTML, que contiene etiquetas procedimentales, como la B de bold , con otras descriptivas como es BLOCKQUOTE y atributo HREF.


(Los ejemplos representan cómo se suele ver el lenguaje de marcado)

\subsection{Marcado de presentación}\label{marcado-presentacion}

Indica el formato del texto, y como su propio nombre indica es utilizado para cambiar la presentación de un texto o documento. En este tipo las etiquetas suelen estar ocultas al usuario.
\\
\\
Ejemplo: Microsoft Word
\\
\\
Cuántas veces habremos oído esa frase...
\\
En muchas ocasiones nos encontramos en sitios nuevos

\subsection{Marcado de procedimientos}\label{marcado-procedimientos}

El lenguaje de marcado de procedimientos también es utilizado para la presentación del texto, pero en este caso las etiquetas sí que son visibles para los usuarios que editan el texto, y se realiza un procesamiento dependiendo de la etiqueta asociada.
\\
\\
Ejemplo: LaTeX
\\
\\
\textbackslash documentclass$[$10pt,a4paper$]$\{article\}
\\
\textbackslash usepackage$[$latin1$]$\{inputenc\}
\\
\textbackslash author\{Juan Pedro Pascual Vitores\}
\\
\textbackslash title\{Introducción\}
\\
\textbackslash begin\{document\}
\\
\textbackslash section\{Introducción\}\textbackslash label{introduccion}
\\
\textbackslash texttt\{``No se dónde está"\}
\\
Cuántas veces habremos oído esa frase...
\\
En muchas ocasiones nos encontramos en sitios nuevos
\\
\textbackslash end\{document\}
\\
\\

\subsection{Marcado descriptivo}\label{marcado-descriptivo}

Se utilizan etiquetas o marcas para describir el texto, sin especificar cómo deben ser presentados, ni en qué orden.

Este tipo de marcado tiene más usos que lo que aparenta a primera vista. Pueden utilizarse para el almacenamiento de metadatos, comunicación de datos (cualquier aplicación pude escribir un documento de texto plano con datos en formato XML), migración de datos (si necesitamos migrar los datos de una base de datos a otra, si ambas trabajan en xml, el trabajo es más sencillo), almacenamiento de gráficos vectoriales (VML) \footnote{VML \textit{Vector Markup Language es un lenguaje XML destinado a la creación de los gráficos vectoriales en 2D o 3D}\cite{vectorml}}, fórmulas matemáticas con XML (MathML)\footnote{MathML \textit{Mathematical Markup Language lenguaje de marcado basado en XML, cuyo objetivo es expresar notación matemática para que distintas máquinas puedan entenderla} \cite{xmlusos} \cite{mathmlwiki}}, Estructuras moleculares e información científica y química con XML ( CML ) \footnote{Chemical Markup Lenguaje\textit{es un lenguaje de marcas basado en XML cuyo objeto es expresar información molecular} \cite{chemicalml}} \cite{desarrollowebxml}
\\
\\
Ejemplo: XML
\\
\\
$<$?xml version=``1.0" encoding=``UTF-8" ?$>$
\\
$<$Documento$>$
\\
	 \tab $<$Texto$>$
	\\
		 \tab \tab $<$Parrafo$>$
		\\
			\tab\tab\tab Cuántas veces habremos oído esa frase...
			\\
			\tab\tab\tab En muchas ocasiones nos encontramos en sitios
			\\
		 \tab \tab $<$\textbackslash Parrafo$>$
		\\
	 \tab $<$\textbackslash Texto$>$
	\\
$<$\textbackslash Documento $>$


\subsection{Lenguajes especializados}\label{lenguajes-especializados}


Los lenguajes especializados podríamos decir que derivan de los lenguajes de marcado descriptivo, ya que se diferencian con el tipo anterior en que se especializan en una materia como podría ser la matemática ( OpenMath o MathML ) o los gráficos (SVG o VRML) \cite{openmath} \cite{xmlusos}.
\\
\\
Ejemplo: OpenMath
\\
\\
$<$OMOBJ xmlns=``url"$>$
\\
 \tab $<$OMA cdbase=``url"$>$
 \\
    \tab\tab$<$OMS cd=``relation1" name=``eq"/$>$
    \\
   \tab\tab$<$OMV name=``x"/$>$
   \\
  \tab$<$OMA\textbackslash $>$
  \\
$<$\textbackslash OMOBJ$>$
\\
Ejemplo: SVG
\\
$<$svg width=``100" height="100"$>$
\\
  \tab$<$circle cx=``50" cy="50" r="40"\textbackslash $>$
  \\
$<$\textbackslash svg4$>$
\\

\section{Algoritmo de Dijkstra}\label{Dijkstra}
En Geoindoor es fundamental el algoritmo de Dijkstra ya que es utilizado para el trazado de rutas a la hora de elegir el camino.
A continuación explicaremos en qué consiste este algoritmo.

El algoritmo de Dijkstra o algoritmo de caminos mínimos es un algoritmo para encontrar el camino más corto en un grafo, desde un punto de partida u origen, hasta uno de los demás puntos.

El grafo consta de vértices y aristas donde cada arista tiene un peso, mediante el cual se calcula cuál es el camino mas corto de un nodo o vértice hasta otro nodo. El principal problema de este algoritmo reside en que es una especialización de la búsqueda de costo uniforme, y no funciona en grafos con aristas de coste o peso negativo, es decir, al elegir siempre el nodo con un peso menor, pueden quedar excluidos nodos que en próximas iteraciones bajarían el costo general del camino \cite{dijkstrabib}. 
\\
\imagenResize{1}{img/Dijkstra}{Ejemplo de grafo con aplicación del algoritmo Dijkstra \cite{dijkstrabib}}{dijkstraimagen}
\\
\subsection{Algoritmo \cite{dijkstrabib}}\label{Explicacion-algoritmo-Dijkstra}
\begin{itemize}
\item
Teniendo un grafo dirigido ponderado de N nodos no aislados.
\item 
Siendo x el nodo inicial.
\item 
Teniendo un vector D de tamaño N en el cual se guardará las distancias desde x al resto de los nodos.
\end{itemize}

\begin{enumerate}
\item
Inicializar la distancia de x en D a 0, debido a que la distancia de x a x es 0. Las demás distancias de D deben ser inicializadas a un valor infinito, ya que al principio son desconocidas.
\item
Siendo a = x tomamos a como nodo inicial.
\item
Recorreremos todos los nodos adyacentes de a, menos los nodos que ya han sido marcados. Llamaremos a los nodos que no han sido marcados $v_i$.
\item
Para el nodo actual, calculamos la distancia desde dicho nodo a los vecinos con la fórmula: dt($v_i$) = Da + d(a,$v_i$), y Si dt($v_i$) $<$ D$v_i$ entonces D$v_i$ = dt($v_i$)
\item
Marcamos el nodo a como completo.
\item
Tomamos como próximo nodo actual el de menor valor en D, y volvemos al paso 3 mientras existan nodos no marcados. También se podría hacer almacenando los valores en una cola de prioridad.

\end{enumerate}

Cuando termine el algoritmo, D contendrá las distancias desde x al resto de los nodos.
\\
La complejidad de este algoritmo es O($|V|^2+|A|$) = O($|V|^2$) sin utilizar cola de prioridad, O(($|A|+|V|$) log $|V|$) = O($|A|$ log $|V|$) utilizando cola de prioridad.

\begin{itemize}
\item
En cada iteración se añade un vértice al subconjunto, para operar y el algoritmo realiza en n-1 iteraciones como máximo.
\item
Se identifica el vértice con la menor etiqueta entre los que no están en el subconjunto, en cada iteración.
\item
Se realiza una operación de suma y una comparación para actualizar la etiqueta de cada uno de los vértices que no están en el subconjunto.

\end{itemize}

Por lo tanto en cada iteración se realiza como máximo 2(n-1) operaciones. Entonces el algoritmo de Dijkstra realiza O($n^2$) operaciones para determinar la longitud del camino más corto entre dos vértices de un grafo. (El grafo debe ser ponderado simple, conexo y no dirigido con n vértices determinados)



\capitulo{4}{Técnicas y herramientas}



\tablaSmall{Manejadores de paquetes y Corredeores de tareas utilizados en el proyecto}{|l|}{manejadores}
{ Manejadores de paquetes y Corredeores de tareas\\ }{ 
Bower \\
Grunt \\
}

\tablaSmall{Formato de intercambio de datos}{|l|}{intercambiodedatos}
{ Formato de intercambio de datos\\ }{ 
JSON\\
}

\tablaSmall{Lenguajes y frameworks utilizados en el proyecto}{|l|}{lenguajes}
{ Lenguajes y Frameworks\\ }{ 
JavaScript \\
Python\\
AngularJS\\
Java (Selenium)\\
Node.js\\
jQuery\\
}

\tablaSmall{Herramientas utilizadas para realizar los test del proyecto}{|l|}{pruebas}
{ Herramientas para pruebas \\ }{ 
Selenium \\
Webserver Stress Tool \\
}

\tablaSmall{Software de control de versiones}{|l|}{controlversion}
{ Control de versiones\\ }{ 
Git\\
ZenHub\\
}

\tablaSmall{Documentación}{|l|}{documentacion}
{ Documentación\\ }{ 
\TeX{}\\
\LaTeX \\
Zotero \\
}

\tablaSmall{Servicios web}{|l|}{servicios}
{ Servicios web\\ }{ 
Heroku\\
Firebase\\
}

\tablaSmall{Herramientas utilizadas para la interfaz del proyecto}{|l|}{interfaz}
{ Interfaz \\ }{ 
HTML5\\
CSS3\\
Bootstrap\\
}

\section{Bower}\label{Bower}

Bower es un gestor de paquetes para la web, es indispensable, ya que nos ha servido para manejar los paquetes del proyecto de una forma simple y rápida. Nos ha permitido manejar las dependencias (angular, bootstrap, jquery, etc\ldots) de forma sencilla y ordenada. 

\section{Grunt}\label{Grunt}

Grunt ha sido de gran utilidad, ya que nos ha permitido compilar el proyecto de forma sencilla y rápida, ayudándonos a tener siempre centralizados los scripts y sus dependencias (imágenes, css, etc\ldots ). A su vez nos dejaba hacer cambios en los scripts en caliente y se volvía a recompilar el proyecto para poder probarlo \textit{in situ}.

Grunt es un corredor de tareas de javascript, que es utilizado para centralizar tareas y scripts, para tener que evitar procedimientos repetitivos. Si se une con Bower permite trabajar con una gran cantidad de paquetes y scripts de forma sencilla.

\section{JSON}\label{JSON}
JSON tiene un papel fundamental en el proyecto ya que al trabajar con JS se hace indispensable, además se ha utilizado en la base de datos.

JSON (JavaScript Object Notation) es un formato de intercambio de datos. Es simple tanto a  la  hora de leerlo, como de escribirlo, es ideal para la transmisión de objetos por la red, y es esa la razón por la que es utilizado.

\section{JavaScript}\label{JavaScript}

JavaScript se considera el lenguaje principal del proyecto, es un lenguaje de scripting, ampliamente utilizado en todos los ámbitos de este proyecto, ya sea para cambiar el flujo de salida, recogida de datos, como para procesamiento de información y llamadas a un sevidor. Recordar que es la base de los frameworks jQuery y AngularJS y para el entorno de ejecución Node.js. Ha sido elegido por su versatilidad y capacidad de adaptarse a cualquier entorno.

\section{Python}\label{python}

Python es un lenguaje de programación interpretado que busca ser un código legible. Es un lenguaje de programación multiparadigma, programación imperativa y programación funcional.

En este proyecto es indispensable ya que es utilizado a la hora de lanzar la aplicación, permitiendo recrear un servidor de forma local.

\section{AngularJS}\label{AngularJS}

AngularJS (Google) es un framework JavaScript de desarrollo de aplicaciones web en el lado cliente, y utiliza el patrón MVC (Model-View-Controller).

AngularJS adquiere las buenas características que tiene JavaScript, pero añade que permite procesar datos de forma mas sencilla y realizar un gran control sobre el flujo de salida.

Este framework se ha utilizado enormemente en este proyecto, ha sido fundamental ya que con ello hemos creado los controladores que permiten interactuar con la información y los usuarios, haciendo control sobre accesos y lo que se permite editar o no.

\section{Java}\label{java}

Java es un lenguaje de programación orientado a objetos, concurrente y de propósito general, y ha sido utilizado en el proyecto en las pruebas, con el WebDriver \footnote{Selenium WebDriver: Colección de enlaces (bindings) específicos para manejar el navegador} y entorno de pruebas de Selenium.

\section{Node.js}\label{Node.js}

Node es un intérprete Javascript en el lado del servidor. Permite construir aplicaciones muy escalables y  maneja miles de conexiones simultáneas en una sólo una máquina física.

En el proyecto ha sido utilizado para crear la Rest API para manejar la base de datos, se ha elegido por su simpleza, aunque cabe destacar que debido a que te permite manejar muchas opciones, es fácil equivocarse y realizar acciones que no se tenían pensadas.

\section{jQuery}\label{jquery}

jQuery podría se considerado una librería más que un framework, esta basada en JavaScript, y su función es simplificar la manera de interactuar con los documentos HTML, con los elementos DOM, manejar eventos, crear animaciones y permite el uso de AJAX.

En este proyecto se ha utilizado tanto sus funcionalidades de AJAX como las manejar eventos y manejar elementos DOM.

\section{Selenium}\label{selenium}

Selenium es un automatizador los navegadores, lo que quiere decir es que permite automatizar pruebas (acciones) directamente sobre una navegador, a través de un WebDriver del propio navegador.

Existen varias librerías para diferentes lenguajes para crear pruebas automatizadas. Pero para este proyecto se ha elegido en lenguaje Java.

\section{Webserver Stress Tool}\label{WebserverStressTool} 

Webserver Stress Tool es una aplicación de prueba HTTP-client/server diseñada para detectar problemas de rendimiento en un sitio web o servidor web.

En este proyecto se ha utilizado para realizar pruebas de rendimiento y estrés.

\section{Git}\label{Git}

Git es un software de control de versiones diseñado por Linus Torvalds, se ha elegido porque es usado por el servicio web heroku así como por GitHub, que es la plataforma en la que se aloja el repositorio del proyecto.

\section{ZenHub}\label{zenhub}

ZenHub es un plug in que permite añadir a GitHub características para mejorar la administración ágil del proyecto, así como mostrar estadísticas sobre el trabajo aplicado al proyecto.

\imagenResize{0.33}{img/zenhubboard}{Tablero ofrecido por ZenHub en GitGub}{zenhubboard}

%\ref{zenhubboard}

\section{\TeX{} y \LaTeX}\label{Latex}

\TeX{} es una herramienta para el manejo de \LaTeX{}, y \LaTeX{} es un sistema para la creación de documentos de texto, de gran calidad tipográfica.
Han sido utilizados para documentar ya que, a pesar de que al principio se hace un poco difícil trabajar con ello, los resultados son excelentes.

\section{Zotero}\label{zotero}

Zotero es una herramienta y plug-in de Firefox que se ha utilizado para la administración de las referencias bibliográficas, nos permite almacenar referencias a las páginas que se visitan a la hora de buscar información por la web. 

\section{Heroku}\label{heroku}

Heroku es una plataforma de servicio de computación en la nube que soporta distintos lenguajes de programación.
En el proyecto ha sido utilizado para sostener el servidor con Node.js con la Rest API. Se ha elegido Heroku por recomendación del tutor, así como sus facilidades para comenzar en ello, ya que ofrece servicios gratuitos y muy buenos tutoriales.

\section{Firebase}\label{firebase}

Firebase es una plataforma de desarrollo de aplicaciones móviles y web, es un servicio de google y nos ofrece hosting, base de datos, almacenamiento, entre otras opciones. Ha sido elegido en este proyecto por su base da datos (JSON) no SQL en tiempo real, que permite trabajar con la base de datos de forma muy rápida y sencilla. El utilizar una base de datos JSON como AnyPlace (aplicación en la que se basa Geoindoor), si se quisiera aunar bases de datos, se conseguiría de una forma relativamente sencilla.


\section{HTML5}\label{HTML5}

HTML5 es un lenguaje de marcado que se ha utilizado para mostrar al usuario la información así como para aprovechar algunas de sus caractrísticas para el procesamiento de información.

Se ha usado en el proyecto porque es muy utilizado, y es soportado por todos los dispositivos.

\section{CSS3}\label{CSS3}

CSS3 lenguaje de hojas de estilo en cascada utilizado para dar una impresión diferente al código HTML, hay que decir que permite crear estilos de forma rápida y sencilla, pero al aumentar el número de clases y el anidamiento es muy difícil de mantener. 
 
\section{Bootstrap}\label{bootstrap}

Bootstrap es un framework para diseño de paginas web y aplicaciones web. Contiene elementos de diseño basados en HTML, CSS y JavaScript. Se ha utilizado para crear la interfaz y aplicarla estilos. 


\capitulo{5}{Aspectos relevantes del desarrollo del proyecto}

En el desarrollo de este proyecto han surgido varios obstáculos que han sido resueltos de diversas maneras. A continuación los explicaremos.

\section{Análisis e investigación de la aplicación }\label{analisis-investigacion}

El principal problema que me encontré fue entender la aplicación sobre la que trabajar y como hacerlo, el principal problema que tuve es que no había una verdadera documentación donde se explicara como funciona al aplicación y cuales son sus funciones y métodos. Con lo cual tuve que investigar y analizar la aplicación a conciencia hasta entenderla.

Otro gran problema que surgió fue entender y crear el servidor que ofrecían, hasta que se decidió trabajar con la aplicación por separado, es decir Architect por un lado y Viewer por otro y después unificarlos para crear el servidor. 

\section{Viewer}\label{viewer}

El problema que surgió con ello, fue a  la hora de crear la función que recoge la posición del usuario para compararla con los puntos determinados. El problema surgió porque en la función getPointLocationUser, watchPosition además de darte las coordenadas del usuario en le momento también te bloquea la la muestra de la posicion del usuario (el punto azul que indica donde esta el usuario), eso es debido a que watchPosition tiene una función callback que bloqueaba. Se consiguió resolver usando clearWatch, para evitar el callback.

\section{Heroku}\label{Heroku}

El principal problema que tuve con ello es que no lo conocía, por lo que tuve que adaptarme y aprender su funcionamiento, aunque hay que decir que los primeros pasos los facilitan gracias a sus tutoriales.

Una vez que se aprendió como funcionaba heroku se comenzo a realizar la Rest API cuyo principal problema fue entender como funcionaba Node.js y las llamadas POST Y GET en Node.js con express.

También es importante decir que con las llamadas al servidor surgieron problemas CORS (access-control-allow-origin),
que se ha solucionado añadiendo las cabeceras de Access-Control-Allow-Origin. 

\section{Firebase}\label{Firebase}

Lo que sucedió con Firebase esta relacionado con su base de datos JSON, cuya estructura de base de datos tuve que idear para que se pudiese conseguir una gran cantidad de información de calidad evitando cantidad de información de poco valor.

La idea es utilizar keys que iguales y añadir una palabra clave para utilizar esa clave como otra key. Es decir la key id2 es la key de un edificio y la key id2rutas tiene las rutas del edificio de key id2.


\section{Architect}\label{Architect} 

El principal problema en la parte de la aplicación de creación de caminos,edificios, y pois, es que las llamadas a al servidor de AnyPlace de determinas url daba problemas de CORS, lo que se soluciono obteniendo información a través de AngularJS 


\section{Otros problemas}\label{otros}

Existieron otros muchos problemas pero que no son comentables ya que no son de tanto interés como estos.

Pero si hay que comentar alguno problema más que me he encontrado, ha sido el uso de jquery para las llamadas, sobre todo con el dataType, el hosting de firebase que no permite hacer llamadas fuera de su dominio y los dynos de heroku.
\capitulo{6}{Trabajos relacionados}

Los trabajos que estan relacionados con este son \href{https://www.google.es/maps}{Google Maps}, \href{https://anyplace.cs.ucy.ac.cy/}{AnyPlace} y \href{https://www.goindoor.co/}{goindoor}.

\section{Google Maps}\label{GoogleMaps}

Google Maps es la semilla del proyecto ya que AnyPlace se basa en ello.
Google Maps ofrece un API en JavaScript que ha sido ampliamente utilizada en este proyecto. La principal ventaja que tiene utilizar el API de Google Maps es el soporte y la atención, además la comunidad que trabaja con Google Maps hace que se resuelvan la dudas fácilmente. 

Hay que reconocer que la documentación que ofrece Google Maps ha sido de gran ayuda, así como su sistema de dibujado sobre el plano. 

La principal diferencia ente Google Maps y Geoindoor es que Google Maps busca la geolocalización en extensiones amplias de terreno, mientras que Geoindoor la busca en interiores, aunque sigue marcando rutas en largas distancias también. Otra diferencia es que Google Maps se centra en crear rutas automáticas para el usuario, mientras que Geoindoor busca crear rutas predeterminadas creadas por usuarios para usuarios.


\section{AnyPlace}\label{AnyPlace}

AnyPlace es la aplicación en la que se basa Geoindoor. AnyPlace siguiendo un poco la línea de Google Maps, busca automatizar las rutas que uno debe seguir, mientras que Geoindoor ofrece rutas predeterminadas creadas por un usuario administrador para que las siga otro usuario cliente; de esta manera se mejora la experiencia de la geolocalización en interiores y permite a los usuarios administradores que los usuarios clientes visiten o localicen los lugares de una forma determinada.

\section{Goindoor}\label{goindoor}

Goindoor es otra sistema utilizado para la geolocalización indoor. La principal diferencia entre Geoindoor y goindoor es que, éste último utiliza bacons o anclajes que son utilizados de sensores para obtener la geolocalización. Mientras que Geoindoor utiliza el sistema de geolocalización de Google Maps y añade rutas predeterminadas, Goindor no tiene tal sistema.

\capitulo{7}{Conclusiones y Líneas de trabajo futuras}

Para concluir puedo decir que el desarrollo ha sido difícil pero satisfactorio en su resultado. La investigación y entendimiento del sistema ha sido dura y ha hecho que la curva de aprendizaje fuera más longeva de lo esperado en un primer momento.

A primera vista se podría pensar que la implementación de funcionalidades es sencilla, pero no es así, esa primera impresión de sencillez desaparece una vez estás desarrollando. 

En cuanto a los resultados, se ha conseguido integrar el agregado de rutas y su interacción.

Para el futuro propondría una serie de mejoras como podría ser: la integración de la base de datos de AnyPlace en Geoindoor, o la unión de ambas herramientas. 

La unión de ambas herramientas sería una gran mejora para ambas partes ya que AnyPlace podría ofrecer rutas predeterminadas y Geoindoor podría llegar a mas gente. Además se evitarían problemas de CORS.

Este proyecto tiene un gran potencial y se podría utilizar en diversas situaciones, como la búsqueda de oficinas en infraestructuras públicas, rutas interactivas en centros comerciales, museos etc..  




\bibliographystyle{plain}
\bibliography{bibliografia}



\end{document}
